\chapter{Literature review}\label{ch:lit}

\section{Epidemic data and policy response}
Our World in Data standardises confirmed cases, deaths, tests and vaccinations across countries \parencite{owid2024,hasell2020,mathieu2021}. Confirmed counts depend on testing effort, which was constrained in Bangladesh, and a high test-positivity rate indicates that many infections went undetected \parencite{hasell2020}.

The Oxford COVID-19 Government Response Tracker (OxCGRT) codes closures, movement restrictions and public-health measures and combines them into a 0--100 stringency index, with a separate economic-support index \parencite{hale2021}. It records announced rules. Behaviour is observed in Google's Community Mobility Reports, which give the percentage change in visits to six categories of place relative to a fixed baseline period in early 2020 \parencite{google2022}. Weekly series like these are serially correlated, and a levels regression overstates its precision unless that is handled \parencite{newey1987}.

\section{Livelihoods, food security and diet}
Early analyses expected the pandemic to threaten food security through lost income \parencite{laborde2020}, as the first evidence from rural Bangladesh suggested \parencite{hamadani2020}.

Women's diet quality is commonly summarised by Minimum Dietary Diversity for Women (MDD-W), the proportion of women aged 15--49 who ate at least five of ten defined food groups in the previous 24 hours, a validated proxy for micronutrient adequacy \parencite{fao2021}. The panel re-analysed here interviewed women in person in November 2019 and by telephone in November 2020 \parencite{panel2021}. A within-person comparison of this kind removes stable differences between women and holds the season fixed. It cannot separate a pandemic-year change from an effect of the change in interview mode, and an average change says nothing about whether poorer and richer households shared it.

\section{Mental health during the pandemic}
Early systematic reviews found elevated rates of anxiety, depression and psychological distress in general populations \parencite{xiong2020,hossain2020}, and the Global Burden of Disease collaborators estimated that major depressive and anxiety disorders became more prevalent in 2020, most where infection rates were highest and mobility fell most \parencite{santomauro2021}.

Bangladeshi studies report high symptom levels among university students \parencite{islam2020} and in a nationwide online survey of adults \parencite{banna2022}. The authors of the survey re-analysed here reported prevalences of 62.5\% for anxiety and 45.3\% for depression, the latter measured by the Patient Health Questionnaire-9 (PHQ-9) at the conventional cut-off of 10 \parencite{kroenke2001,manea2012}. Using hierarchical linear regression, they attributed 5--23\% of the variance in outcomes to factors they called non-modifiable (student status, female sex, being unmarried, northern residence and rural residence) and a further 10--25\% to factors they called modifiable (mindfulness, optimism and vaccination) \parencite{islam2023}.

Online convenience samples over-represent young, urban and educated respondents, so a prevalence figure describes the sample. An unsupervised online form admits duplicated or templated entries, which a regression will absorb without warning. The predictors labelled modifiable are self-reports collected in the same questionnaire as the outcomes, so their association with symptoms is concurrent and its direction is not identified.

\section{Prediction models and latent structure}
The learners compared in this report are penalised logistic regression \parencite{tibshirani1996,zou2005}, random forests \parencite{breiman2001} and gradient-boosted trees \parencite{friedman2001,chen2016}. Reporting guidance for clinical prediction models requires performance to be estimated on data not used for fitting or tuning \parencite{collins2015}, and calibration, the agreement between predicted risks and observed frequencies, to be reported alongside discrimination \parencite{vancalster2019}. A systematic review of 71 clinical prediction studies found no performance benefit of machine learning over logistic regression in comparisons at low risk of bias \parencite{christodoulou2019}. The choice of predictors affects what a measure of discrimination such as the area under the ROC curve (AUC) means. Scores from scales administered together with the outcome are concurrent correlates, and a model that uses them measures shared information without showing that cases could be identified in advance.

Latent profile analysis fits a Gaussian mixture with diagonal covariance to continuous indicators and chooses the number of profiles by the Bayesian information criterion (BIC) or the integrated completed likelihood (ICL) \parencite{scrucca2016,biernacki2000}. Profiles fitted to correlated symptom scales can differ in level and hardly at all in shape, in which case they are bands along one severity dimension. A mixture likelihood also rewards a narrow component placed on identical observations, so duplicated records can surface as a class of their own.
